\section{Discussion, Limitations, Future Work}
In this work, we introduced \emph{insight anticipation} as a measurable paradigm for automated scientific discovery. By isolating the insight generation phase, we demonstrated that models can effectively synthesize the core insight of a downstream paper when provided with its foundational parent papers. Our results using \ourmethod{} indicate that the trajectory of scientific intuition is partially predictable, and that optimizing language models via reinforcement learning with similarity-based rewards is a highly effective training strategy for this task. However, this foundational formulation leaves several important limitations and open questions.

One limitation of the work is that we assume that downstream contributions derive from two parent papers due to context constraints, despite research ideas being heavily shaped by broader intellectual contexts. Furthermore, parent identification is imperfect. Citations do not always reflect true conceptual influence, and influential ideas may remain uncited. Finally, by assuming an oracle literature selection criterion, we explicitly decoupled insight generation from parent selection, which may not be feasible for some scientific tasks.

Future research can integrate automated retrieval systems to unify parent selection and synthesis within a single end-to-end framework. Additionally, extending the task to accommodate multi-source lineage and developing evaluation metrics that prioritize conceptual novelty over textual similarity could be interesting. Ultimately, testing these models in active, human-in-the-loop research settings will determine their true potential as catalysts for scientific discovery.



% First, to enable controlled evaluation within \ourbenchmark{}, our abstractions assume that downstream contributions derive from exactly two parent papers. While this pairwise framing successfully isolates the mechanics of targeted synthesis, real-world scientific breakthroughs are rarely so cleanly decomposable. Research ideas are heavily shaped by broader intellectual contexts, encompassing multiple prior works, tacit domain knowledge, and evolving community trends. Constricting the lineage to two parents underestimates the complexity of true ideation and may bias models toward overly simplistic, pairwise reasoning patterns.

% Even within our constrained setup, parent identification is imperfect. Citations do not always reflect true conceptual influence, and influential ideas may remain uncited. Conversely, some citations serve rhetorical or historical purposes rather than indicating substantive intellectual dependence. As a result, the mapping from cited parents to “true” idea sources is noisy. This noise places an upper bound on achievable performance and complicates interpretation of both model errors and evaluation scores.

% We explicitly isolate insight synthesis from parent selection by assuming that relevant parents are given. However, selecting which prior work to build upon is itself a central component of scientific reasoning. 

% Our approach has several limitations:
% \begin{itemize}
%     \item Real papers often draw from many influences, whereas we restrict attention to two parents.
%     \item Not all influential ideas are explicitly cited, introducing noise into parent identification.
%     \item Parent selection is assumed to be given and is not modeled.
% \end{itemize}



% This work provides evidence that scientific insight synthesis is partially predictable and that reinforcement learning is particularly well-suited for modeling this process.

% Future directions include scaling to more parent papers, jointly modeling parent selection and idea synthesis, and studying whether models can anticipate emerging research directions rather than reconstructing known ones.